\documentclass{exam}
\usepackage{../commonheader}

\discnumber{2}
\title{Python Functions and Expressions}
\date{September 2 -- September 5, 2025}

\begin{document}
\maketitle

\begin{meta}
    \textbf{Example Timeline}
    \begin{itemize}
        \item Teaching tips:
            \subitem{\textbf{Each problem's "Suggested Time" is not assigned with the 1 hour time for sections in mind.} It is more of a note of how long each individual problem may take. That is, if you added all the "Suggested Times" in the coding section, it could very well be greater than 1 hour (especially in more complicated worksheets in the future).}
            \subitem{\textbf{However, the time in brackets will add up to ~50 minutes, and is intended to be a suggested outline for your section.}}
        
        \item Intros, Ice Breakers, \& Logstics [10 min]
            \subitem{Introduce yourself, and give an overview of what CSM offers. Start off with an icebreaker of your choice!}
        \item Mini-lecture: Python Expressions and Functions [10 min]
        \item Q1: Potpourri [7 min]
            \subitem{You can skip the question with the if statements, and feel free to skip some parts!}
        \item Q2: Identify Types [7 min]
        \item Q3: Funky Functions [7 min]
        \item Q4: Fahrenheit to Celsius [5 min]
        \item Q5: Score Needed [7 min]
        \item Q6: WWPD [If Extra Time]
            
    \end{itemize}
\end{meta}

\section{Expressions}
\begin{questions}
    \subimport{../../topics/basics/}{potpurri.tex}
    \begin{questionmeta}
        Suggested Time: 7 min
        \begin{itemize}
            \item You can probably skip a lot of these as soon as your students get the idea.
            \item Your students might be new to solving "compound" questions like the last one; Consider giving them tips on how to break down such problems.
        \end{itemize}
    \end{questionmeta}

    \subimport{../../topics/functions-and-expressions/medium/}{identify-types.tex}
    \begin{questionmeta}
        Suggested Time: 7 min, Difficulty: Medium
        \begin{itemize}
            \item Make sure students understand why '3' + 4 errors, as well as how string concatention works in the case of '3' + '4'
        \end{itemize}
    \end{questionmeta}
\end{questions}


\section{Functions}
\begin{questions}
    \subimport{../../topics/basics/}{funky-funcs.tex}
    \begin{questionmeta}
        Suggested Time: 7 min
        \begin{itemize}
            \item Show students how to tackle problems with nested expressions (like the process in the solution for this problem).
        \end{itemize}
    \end{questionmeta}

    \subimport{../../topics/functions-and-expressions/easy/}{celsius.tex}
    \begin{questionmeta}
        Suggested Time: 5 min, Difficulty: Easy
    \end{questionmeta}

    \subimport{../../topics/functions-and-expressions/medium/}{score-needed.tex}
    \begin{questionmeta}
    Suggested Time: 7 min, Difficulty: Medium
    Make sure that students understand the question prompt and how to come up with the formula.
    \end{questionmeta}

    \subimport{../../topics/functions-and-expressions/medium/}{fgh.tex}
    \begin{questionmeta}
        Suggested Time: 5 min, Difficulty: Medium
        \begin{itemize}
            \item Optional, for if you have time at the end!
        \end{itemize}
    \end{questionmeta}
\end{questions}

\end{document}